\section*{Chapter \ref{ch:dv:convertibles-and-the-credit-equity-link} --- Convertibles and the Credit-Equity Link}

\begin{solution}{exo:dv:convertibles-and-the-credit-equity-link:1}
\omterm{def:dv:convertibles-and-the-credit-equity-link:price}{Conversion price} $100/2.5=40$; \omterm{def:dv:convertibles-and-the-credit-equity-link:value}{conversion value} $2.5\times36=90$; premium $110/90-1=22.2\%$.
\end{solution}

\begin{solution}{exo:dv:convertibles-and-the-credit-equity-link:2}
The holder can always keep the bond and receive its coupons and principal (worth the \omterm{def:dv:convertibles-and-the-credit-equity-link:floor}{bond floor}), or convert now (worth the \omterm{def:dv:convertibles-and-the-credit-equity-link:value}{conversion value});
a right to choose later is worth at least each of the two, so the convertible is worth at least their maximum.
\end{solution}

\begin{solution}{exo:dv:convertibles-and-the-credit-equity-link:3}
When the \omterm{def:dv:convertibles-and-the-credit-equity-link:value}{conversion value} exceeds the call price, the issuer can force conversion and so remove the holder's remaining time value; the holder is
short that call. The effect is largest where the time value is still large and the call can be exercised: near and above the trigger, here 4.7 at
52 against 3.9 at 40.
\end{solution}

\begin{solution}{exo:dv:convertibles-and-the-credit-equity-link:4}
Under the pricing measure the share must return $r-q$ on average. With intensity $\lambda(S)$ it jumps to zero, an expected loss of $\lambda(S)$ per
unit time; the diffusion's drift must add it back: $\E[dS/S]=(r-q+\lambda)\,dt-\lambda\,dt=(r-q)\,dt$.
\end{solution}

\begin{solution}{exo:dv:convertibles-and-the-credit-equity-link:5}
A wider spread lowers the \omterm{def:dv:convertibles-and-the-credit-equity-link:floor}{bond floor}, but part of the convertible's value comes from the conversion right, which is worth more than the floor's
loss in many scenarios (and gains from the higher share drift in the model); only the bond-like part carries the full sensitivity: 0.0167 per bp
against 0.0296.
\end{solution}

\begin{solution}{exo:dv:convertibles-and-the-credit-equity-link:6}
In the \omterm{def:dv:convertibles-and-the-credit-equity-link:e2c}{equity-to-credit model} a falling share raises the hazard and lowers the \omterm{def:dv:convertibles-and-the-credit-equity-link:floor}{bond floor}, so the convertible falls faster with the share: the
delta adds a credit channel (1.39 against 0.84 at 20). Hedging with the constant-hazard delta leaves the desk under-hedged in exactly the falls
that widen the credit, implicitly long credit.
\end{solution}

\begin{solution}{exo:dv:convertibles-and-the-credit-equity-link:7}
With $p=2$ the delta at 30 is 1.65 and the delta-hedged loss shrinks to 1.25 per 100 of face (2.55 with $p=1.2$): the stronger link puts more of
the credit move into the equity delta, which the short shares then hedge. The convertible is worth 97.32 today and 86.18 after the move.
\end{solution}

\begin{solution}{exo:dv:convertibles-and-the-credit-equity-link:8}
The hedges are local and model-based: the delta depends on the model of the credit link, the CDS covers small spread moves but not the jump to
default fully, and neither covers the convertible cheapening relative to its model value when financing is withdrawn and holders sell together (2008).
The book is also short the combination of equity and credit moves in distress, where the gamma turns negative.
\end{solution}

\begin{solution}{pb:dv:convertibles-and-the-credit-equity-link:1}
\textbf{1.} 98.47; \omterm{def:dv:convertibles-and-the-credit-equity-link:floor}{bond floor} 82.91; \omterm{def:dv:convertibles-and-the-credit-equity-link:value}{conversion value} 75; premium 31.3\%.
\textbf{2.} 1.55 shares per bond: 1.55 million shares short for 100 million of face (one million bonds).
\textbf{3.} 69.07 at a share price of 10, where the hazard is 18.7\% a year.
\textbf{4.} 4.65 per bond: 138.64 with the call against 143.29 without.
\textbf{5.} Because the hazard, and so the discounting of the bond's cash flows, depends on the share price.
\textbf{6.} 0.0167 per basis point per bond: 16\,700 per basis point for the position.
\textbf{7.} 40.6 million of protection (0.0167 per bond over a risky annuity of 4.12 per unit of spread).
\textbf{8.} 2.21 per 100 of face, 2.21 million: the convexity.
\textbf{9.} 3.35 million.
\textbf{10.} 1.37 shares per bond: 0.18 million fewer shares short.
\textbf{11.} 86.62; a loss of 2.55 million.
\textbf{12.} A loss of 4.38 million.
\textbf{13.} A gain of 1.93 million (the protection gains 4.47 million).
\textbf{14.} The spread widening alone costs 3.35 and the share fall alone earns 2.21 of convexity; the choice of delta is worth 1.83 (4.38 against 2.55).
\textbf{15.} The convertible cheapening relative to its model value as holders sell together (financing withdrawn), and the jump to default.
\textbf{16.} The equity-to-credit delta, calibrated to the issuer's credit curve and options: it hedges the credit channel that the scenario exploits.
\textbf{17.} At least the credit sensitivity at the current price, more if the jump to default at the current recovery is large relative to the book;
not more than the desk is willing to pay for in carry.
\textbf{18.} Limits on the joint equity-down, credit-wider stress, on concentration by issuer, on leverage and on the share of the book financed short-term.
\textbf{19.} 2.55 per 100 of face, 2.55 million on 100 million, with the equity-to-credit delta (4.38 million with the constant-hazard delta).
\textbf{20.} Convexity in the share, bought with credit risk and financed with borrowed money.
\end{solution}

\begin{solution}{iq:dv:convertibles-and-the-credit-equity-link:1}
Left: close to the \omterm{def:dv:convertibles-and-the-credit-equity-link:floor}{bond floor}, which itself falls as the share falls when credit deteriorates (distressed region). Middle: convex, both bond and
option (the hybrid region, where the arbitrage lives). Right: close to the \omterm{def:dv:convertibles-and-the-credit-equity-link:value}{conversion value}, equity-like, capped near the call trigger by the
issuer's call.
\iqlookfor{the three regions and the call's cap.}
\end{solution}

\begin{solution}{iq:dv:convertibles-and-the-credit-equity-link:2}
In $x=\ln S$: $V_t+\frac12\sigma^2V_{xx}+(r-q+\lambda-\frac12\sigma^2)V_x-(r+\lambda)V+\lambda RF=0$. Terminal: $\max(F+\text{coupon},\kappa S)$.
Constraints after each step: $V\ge\kappa S$; $V\le\max(\text{call price},\kappa S)$ when callable; $V\ge$ put price on put dates; coupons added at
coupon dates. Boundaries: recovery at very low $S$, \omterm{def:dv:convertibles-and-the-credit-equity-link:value}{conversion value} at very high $S$.
\iqlookfor{the hazard terms, and the constraints as projections.}
\end{solution}

\begin{solution}{iq:dv:convertibles-and-the-credit-equity-link:3}
The convexity (gamma trading against the short shares), the coupon net of the short's borrow cost, and convergence of cheap bonds to value. Risks:
credit (spread and default), the model's delta, liquidity and financing (forced sales), the short squeeze or borrow recall on the shares, and
issuer calls.
\iqlookfor{sources of return and the risk list.}
\end{solution}

\begin{solution}{iq:dv:convertibles-and-the-credit-equity-link:4}
Step the PDE back in time (Crank--Nicolson, with implicit start-up steps after each discontinuity), then project: take the maximum with the
\omterm{def:dv:convertibles-and-the-credit-equity-link:value}{conversion value}, the minimum with the call price (or conversion) where the call is live and its trigger met, the maximum with the put price on put
dates; add coupons at their dates. Handle soft-call triggers observed over days with an extra state or an approximation, and document it.
\iqlookfor{projection after each step, and care at discontinuities.}
\end{solution}

\begin{solution}{iq:dv:convertibles-and-the-credit-equity-link:5}
Joint scenarios of share and credit: shares down 10--40\% with spreads 100--500 basis points wider at the new price, jumps to default of the largest
names, a volatility shock, a borrow recall, and a liquidity scenario in which convertibles cheapen by several points relative to model with no move in
shares or credit.
\iqlookfor{joint equity--credit moves plus liquidity.}
\end{solution}

\begin{solution}{iq:dv:convertibles-and-the-credit-equity-link:6}
The book must sell about half its positions or find new financing. Sales of the same bonds by many funds push their prices below model values; the
hedges (short shares, CDS) do not offset that cheapening; losses reduce capital and force further sales. Prices recover only as new capital arrives,
slowly.
\iqlookfor{forced sales, cheapening relative to model, and slow recovery.}
\end{solution}
